\scanpage{41}
Théorèmes de finitude pour le groupe fondamental

1) Sur un corps \unsure{séparablement|algébriquement} clos:

\[
\text{a)}\qquad X\ \text{propre}\,/k\quad\Longrightarrow\quad\pi_1(X)\ \text{de type fini}
\]

\diagram{\begin{tikzcd}[ampersand replacement=\&,column sep=small,row sep=large,execute at end picture={\draw[gray] (pb.north) .. controls ([xshift=10mm]pb.north east) and ([xshift=8mm]pc.south east) .. (pc.south) .. controls ([xshift=-8mm]pc.south west) and ([xshift=-8mm]pb.north west) .. (pb.north);}]
{\text{b)}\quad X\ \text{t.f.}\,/k} \& \Longrightarrow \& |[alias=pb]| {\pi'_1(X)} \& {\text{de }\boxed{\text{\uline{présentation}}}\ \text{\uuline{finie}}}\\
{\text{c)}\quad X\ \text{t.f.}\,/k} \& \Longrightarrow \& |[alias=pc]| {\pi_1^t(X)} \arrow[u] \& {\text{de type fini}}
\end{tikzcd}}

\margin{($'$ indique partie première à la caractéristique)}

\margin{($t$ indique que la ramification est tame sur \unsure{tt|H} \deleted{\unsure{modèle} propre} \inserted{``diviseur''} \deleted{d'une \unsure{comp.|corps}} du corps de fonctions des composantes irréd]}

[N.B. c) plus général que a)]

\margin{$X_i\subset X$ irréd. \unsure{composantes|composant}}

a) \uline{démontré}, [par \uline{descente} réduction au cas normal, par \uline{Bertini} au cas d'une courbe, puis théorie de spéc. du gpe fond. et méthode \unsure{transcendante|transcendance}]

\deleted{b) \inserted{\unsure{le type fini est prouvé}} [réduit au cas $X$ normal [descente] \unsure{par}]}

b) Le cas ``\uline{type fini}'' est \unsure{prouvé|présumé} [réduction au cas normal par descente puis au cas non singulier ``joli'' du sens Artin, on utilise alors la suite exacte des $\pi_1$ i.e. th. de spécialisation du \unsure{grpe.|gpe} fond, qui nous ramène au cas d'une courbe normale, où on est ramené par th. de spéc. au cas \unsure{transcendant}]

\unsure{Enfin|A priori|Espoir|Quid} : \unsure{ceci}, on est ramené pour b) au cas $X$ \uline{normal} [descente] et \unsure{même} [via \uuline{résolutions}] au cas $X$ \deleted{\uline{régulier}}\inserted{lisse/$k$}. \uline{Si} $X$ \deleted{\unsure{projectif|propre}]}\unsure{est|on est} [\uline{Lefschetz}] ramené, au cas d'une surface. Si on a \unsure{résolution} sous forme forte, \uline{cas} d'une surface \unsure{projective}, \deleted{intersection} \deleted{\unsure{complète|complété}}\inserted{lisse} diviseur (à normal crossings).

\deleted{c)}
% archive sheet number: 40
